\section{SPACE COMBAT}

Space combat combines \Imperium's range and weapon decisions with the
\SpaceEmpires\ firing order, attack and defense values, screening, and hull
damage. Every combat round occurs at either \textbf{Long Range} or
\textbf{Close Range}. Missiles dominate the approach; beams become decisive
after fleets close.

The two bands deliberately compress Traveller's more detailed scale. Picket
ships and Tactics stand in for thrust, sensors, and fleet handling. Point
defense lasers protect against missile salvos, while sandcaster reactions
protect against direct-fire beams.

\subsection{Armament Packages}

Every SC, DD, CA, BC, BB, and DN group has one Armament Package. Record the
package on its technology row. All ships in a group must have the same package.

\begin{tabularx}{\columnwidth}{@{}lXX@{}}
\toprule
Package & Long Range & Close Range \\
\midrule
Beam (B) & Cannot fire & Beam at full Attack \\
Missile (M) & Missile at full Attack & Missile at $-1$ Attack, after beams \\
Mixed (X) & Missile at $-1$ Attack & Beam at $-1$, or missile at $-2$ after beams \\
\bottomrule
\end{tabularx}

``Attack'' in this table means the printed Attack value before technology,
Fleet Size, defense, and other modifiers. A modified Hit Number of zero or less
still hits on a natural 1.

Fighters and Ship Yards have a fixed Beam Package. Bases have a fixed Missile
Package. Carriers, Transports, and Fleet Trains have no package and cannot
fire if their printed Attack is zero.

Choose a package when a ship is purchased. A ship may change package only
during peace, at a cost of 1 RU per Hull Size. Package identity is hidden with
the group until the group is revealed.

\subsection{Prepare the Battle}

Move the opposing groups off-map and mark the contested system. Reveal all
ships, Armament Packages, and current technology. Arrange each side's groups by
firing class.

Every Base and Ship Yard on a friendly, unneutralized surface in the system
joins its owner's space battle. In a binary system, installations on every
friendly surface may join. Bases and Ship Yards are combat-capable, count for
Fleet Size and screening, and cannot retreat. Ground Units and surfaces do not
otherwise participate in space combat.

The player whose movement caused the battle is the attacker. If neither side
moved into the system this phase, the Phasing Player is the attacker.

\subsection{Initial Range}

The first round begins at Long Range. This represents missile exchanges while
the opposing formations approach. A scenario or card may change the initial
range.

\subsection{Range Control}

Before the second and every later round, each player selects one
combat-capable, non-Fighter ship as a \textbf{Picket}. A Picket still fires
normally. Its Maneuver Rating is:

\begin{tabularx}{\columnwidth}{@{}Xr@{}}
\toprule
Ship & Maneuver \\
\midrule
Scout & +2 \\
Destroyer & +1 \\
Cruiser, Carrier, or Transport & 0 \\
Battlecruiser, Battleship, or Dreadnought & $-1$ \\
No eligible ship & $-2$ \\
\bottomrule
\end{tabularx}

Fleet Trains, Bases, and Ship Yards cannot be Pickets. Each player rolls a d10
and adds the Picket's Maneuver Rating and friendly Tactics. The side with
fewer combat-capable non-Fighter ships also adds +1. The higher total chooses
Long or Close Range for the round. On a tie, range remains unchanged.

This roll replaces \Imperium's range-determination roll and Traveller's
individual thrust accounting. Destroying an enemy Scout can therefore change
control of the engagement.

\subsection{Screening and Fleet Size}

After establishing range each round, count combat-capable ships, including
operational Fighters, Bases, and Ship Yards. The side with more may screen
eligible groups under \SpaceEmpires\ rule 5.7. Fleet Trains and other
non-combat units are screened automatically while a friendly combat-capable
ship remains.

If one side has at least twice as many unscreened combat-capable ships as the
other, it receives the normal Fleet Size Bonus of +1 Attack for that firing
round. Recalculate range control, screening, and Fleet Size every round.

\subsection{Defensive Reactions}

Defensive reactions are declared after screening and before any ship fires.

\subsubsection{Point Defense}

At Point Defense 1 or higher, an unscreened SC or DD with a Beam or Mixed
Package may escort one friendly target group. The escort forfeits all offensive
fire and Close Assaults this round. The protected group gains +1 Defense
against missiles, increased to +2 at Point Defense 2 or 3. A target group may
have only one escort. If the escort is destroyed before a missile attack, the
bonus is lost.

At Point Defense 3, the escort may also make one normal attack at $-1$ Attack.
This is an exception to forfeiting offensive fire. Point defense represents
laser interception, electronic warfare, and terminal countermeasures against
an incoming salvo.

\subsubsection{Disperse Sand}

Any non-Fighter combat ship may disperse sand. It forfeits all offensive fire,
Close Assaults, and point defense this round, but gains +1 Defense against
every beam attack for the round. Sand has no effect against missiles. A ship
that has already fired cannot disperse sand.

\subsection{Firing Order}

At \textbf{Long Range}, resolve missile fire from Class A through E. Beam-only
ships cannot fire.

At \textbf{Close Range}, first resolve all beam fire from Class A through E.
Then resolve short-range missile fire from Class A through E. A ship may fire
only once in the round; a Mixed ship must choose beam or missile fire. A ship
destroyed by beam fire cannot later make a short-range missile attack.

Within the same class and weapon step, higher Tactics fires first. The defender
fires first if Tactics is tied. Fire is not simultaneous. A destroyed ship does
not fire later in the round.

Each surviving firing ship may target any unscreened enemy unit. A group's
ships may choose different targets, and the owner may observe each die before
choosing the next target.

\subsection{Fire Resolution}

Start with the Attack value after its Armament Package modifier, then
calculate:

\begin{center}
\textbf{Hit Number = Attack + Attack Tech + Fleet Bonus \\
$-$ target Defense $-$ target Defense Tech \\
$-$ defensive reaction bonuses}
\end{center}

Roll a d10. A result equal to or less than the Hit Number scores one hit. A
natural 1 always hits. Technology added to a particular ship may not exceed
that ship's Hull Size. Package, card, and combat-option modifiers are applied
after enforcing that technology limit.

Damage is cumulative. When a ship has hits equal to its Hull Size, remove it
and reduce its group's numeral. Carry surviving damage out of combat; it may be
repaired in a later Logistics Phase.

\subsection{Missile Options}

\subsubsection{High-Intensity Missile Fire}

A ship making a missile attack at either range may expend its magazines for
+1 Attack. Mark it Missile Depleted after firing. It cannot make another
missile attack during the current Combat Phase. A Mixed ship may still fire
beams in a later round.

High-Intensity Fire is declared separately for each ship when it fires. A ship
that is retreating or performing point defense cannot use it.

\subsubsection{Short-Range Missile Fire}

At Close Range, Missile ships attack at $-1$ and Mixed ships at $-2$. These
attacks occur only after every beam class has fired. High-Intensity Fire may be
combined with short-range missile fire.

\subsection{Close Assault}

At Close Range, a SC, DD, or Fighter making a beam attack may declare a Close
Assault. Before the attacker rolls, its target may make immediate defensive
beam fire if it has a Beam or Mixed Package and has not fired, retreated, or
used a defensive reaction this round. This defensive shot expends the target's
normal fire.

If the assaulting ship survives, it attacks with +1 Attack. If defensive fire
destroys it, it is removed before it can fire. This replaces \Imperium's
Suicide Attack while retaining its risk-and-reward role.

\subsection{Fighters and Carriers}

Use \SpaceEmpires\ section 11 except as changed here. Fighters use Beam fire
and may conduct Close Assaults; they cannot fire at Long Range. Fighters in a
system require enough friendly Carrier capacity or a friendly unneutralized
World or Outpost. Excess Fighters are destroyed when their final base or
carrying capacity is lost.

\subsection{Landing During a Space Battle}

After the first firing round, either side may land carried Space Marines on a
surface in the system. After the second and each later round, either side may
land any carried Ground Units. The attacker announces first each round.

A Transport may land units even if screened. Space combat then continues.
Ground combat waits until space combat in the system is complete. A player may
therefore win on the ground after losing in orbit, but surviving enemy ships
may attack their Transports before the space battle ends.

\subsection{Retreat}

Beginning in the second firing round, a group may retreat instead of firing
when its printed class is reached in the applicable weapon step. A group with
no usable weapon at the current range may still retreat at its printed class.
Move it to one legal adjacent system. Carriers may not abandon excess Fighters.
Bases, Ship Yards, Fleet Trains, and other non-combat units cannot retreat.

Combat ends when only one side has combat-capable ships in the stellar hex.
Destroy all enemy non-combat ships left with no escort. Remove Missile Depleted
and defensive-reaction markers at the end of the Combat Phase.

\subsection{Repair}

During Maintenance, remove damage from ships at a connected friendly World by
paying 1 RU per hit. Damage at any other location cannot be repaired. Repair is
optional and occurs after normal maintenance; a damaged ship that is scrapped
needs no repair.
